\documentclass[11pt,a4paper]{article}

\usepackage[T1]{fontenc}
\usepackage[utf8]{inputenc}
\usepackage{lmodern}
\usepackage[ngerman]{babel}
\usepackage[a4paper,margin=2.2cm]{geometry}
\usepackage{array}
\usepackage{longtable}
\usepackage{tabularx}
\usepackage{xcolor}
\usepackage{hyperref}

\setlength{\parindent}{0pt}
\setlength{\parskip}{0.5em}
\renewcommand{\arraystretch}{1.35}

\newcommand{\term}[1]{\fcolorbox{black!20}{blue!6}{\strut #1}}
\newcommand{\solution}[1]{\textcolor{black}{#1}}

\title{\textbf{Lösungsblatt: Kategorien von Systemkomponenten}}
\author{Modul 321 -- Lektion 06}
\date{}

\begin{document}

\maketitle

\fcolorbox{blue!40}{blue!6}{
  \parbox{\dimexpr\linewidth-2\fboxsep-2\fboxrule-6pt\relax}{
    \textbf{Hinweis:} Dieses Lösungsblatt enthält Musterlösungen und einen möglichen
    Erwartungshorizont. Andere Antworten sind korrekt, wenn sie fachlich sauber
    begründet sind und die Rollen der Komponenten nachvollziehbar eingeordnet werden.
  }
}

\section*{1. Begriffe erklären}
\begin{longtable}{|>{\raggedright\arraybackslash}p{0.30\linewidth}|>{\raggedright\arraybackslash}p{0.62\linewidth}|}
\hline
\textbf{Begriff} & \textbf{Mögliche Lösung} \\
\hline
\term{Identitätsverwaltung} &
\solution{Die Identitätsverwaltung kümmert sich um Benutzeridentitäten, Anmeldung, Rollen und Berechtigungen. Sie prüft also, wer eine Person ist und worauf sie zugreifen darf.} \\ \hline
\term{Ereignisverwaltung} &
\solution{Ereignisverwaltung verarbeitet oder verteilt Meldungen über eingetretene Zustandsänderungen. Dadurch können mehrere Teile eines Systems auf ein Ereignis reagieren.} \\ \hline
\term{Monitoring} &
\solution{Monitoring beobachtet den Zustand und das Verhalten eines Systems. Es sammelt zum Beispiel Logs, Metriken oder Fehlerraten, damit Probleme sichtbar werden.} \\ \hline
\term{Datenhaltung} &
\solution{Datenhaltung speichert Informationen dauerhaft und stellt sie für Lese- und Schreibzugriffe bereit. Dazu gehören zum Beispiel Datenbanken oder Dateispeicher.} \\ \hline
\term{Lastenverteilung} &
\solution{Lastenverteilung verteilt Anfragen oder Arbeit auf mehrere Instanzen. Dadurch kann ein System stabiler laufen und besser skalieren.} \\ \hline
\end{longtable}

\section*{2. Beispiele zuordnen}
\begin{enumerate}
  \item \solution{Identitätsverwaltung}
  \item \solution{Ereignisverwaltung}
  \item \solution{Monitoring}
  \item \solution{Datenhaltung}
  \item \solution{Lastenverteilung}
  \item \solution{Identitätsverwaltung}
\end{enumerate}

\section*{3. Komponenten vergleichen}
\begin{tabularx}{\linewidth}{|>{\raggedright\arraybackslash}p{0.28\linewidth}|>{\raggedright\arraybackslash}X|>{\raggedright\arraybackslash}X|}
\hline
\textbf{Frage} & \textbf{Monitoring} & \textbf{Datenhaltung} \\
\hline
Welches Hauptziel hat die Komponente? & \solution{Systemzustand beobachten und Probleme sichtbar machen} & \solution{Daten dauerhaft speichern und verfügbar halten} \\ \hline
Welche typischen Informationen verarbeitet sie? & \solution{Logs, Metriken, Fehlerraten, Antwortzeiten} & \solution{Geschäftsdaten, Benutzerdaten, Dateien, Datensätze} \\ \hline
Was wäre ein typischer Fehlerfall? & \solution{Alarme fehlen oder Fehler bleiben unbemerkt} & \solution{Datenverlust, inkonsistente oder nicht verfügbare Daten} \\ \hline
\end{tabularx}

\vspace{1em}

\begin{tabularx}{\linewidth}{|>{\raggedright\arraybackslash}p{0.28\linewidth}|>{\raggedright\arraybackslash}X|>{\raggedright\arraybackslash}X|}
\hline
\textbf{Frage} & \textbf{Identitätsverwaltung} & \textbf{Lastenverteilung} \\
\hline
Welche Kernaufgabe hat die Komponente? & \solution{Authentisierung und Autorisierung} & \solution{Anfragen auf mehrere Instanzen verteilen} \\ \hline
Welche Probleme löst sie typischerweise? & \solution{Wer darf zugreifen, zentrale Rechteverwaltung, sichere Anmeldung} & \solution{Überlastung einzelner Instanzen, bessere Skalierung, Ausfallschutz} \\ \hline
Warum sollte man sie nicht verwechseln? & \solution{Weil Sicherheit und Zugriffssteuerung eine andere Verantwortung sind als Betriebs- und Lastverteilung.} & \solution{Weil ein Load Balancer keine Benutzerrechte prüft und ein Login-System keine Last verteilt.} \\ \hline
\end{tabularx}

\section*{4. Aussagen beurteilen}
\begin{enumerate}
  \item \solution{Falsch. Monitoring sammelt vor allem Betriebsinformationen, nicht primär Geschäftsdaten.}
  \item \solution{Richtig. Eine zentrale Identitätsverwaltung kann von mehreren Anwendungen gemeinsam verwendet werden.}
  \item \solution{Richtig. Lastenverteilung hilft, Anfragen auf mehrere Instanzen zu verteilen und damit Skalierung zu unterstützen.}
  \item \solution{Falsch. Auch in eventbasierten Systemen müssen Daten weiterhin gespeichert werden.}
  \item \solution{Richtig. Eine Datenbank speichert Daten, entscheidet aber nicht automatisch sauber über Rollen und Berechtigungen.}
\end{enumerate}

\section*{5. Lernplattform analysieren}
\begin{enumerate}
  \item \solution{Anmelden: Identitätsverwaltung. Dateien speichern: Datenhaltung. Andere Teile informieren: Ereignisverwaltung. Ausfälle sichtbar machen: Monitoring. Viele Zugriffe abfangen: Lastenverteilung.}
  \item \solution{Besonders sensibel ist die Identitätsverwaltung, weil dort Anmeldung und Berechtigungen geregelt werden.}
  \item \solution{Für Stabilität und Skalierung ist besonders die Lastenverteilung wichtig.}
\end{enumerate}

\section*{6. Fehlende Komponente erkennen}
\begin{enumerate}
  \item \solution{Monitoring und Lastenverteilung.}
  \item \solution{Monitoring fehlt, weil Fehler erst von Benutzenden gemeldet werden. Lastenverteilung fehlt oder ist zu schwach, weil das System bei hoher Last langsam wird.}
  \item \solution{Mögliche Folgen sind schlechte Reaktionszeiten, Ausfälle, schwer erkennbare Fehler und verspätete Problemanalyse.}
\end{enumerate}

\section*{7. Mini-Entwurf}
\begin{enumerate}
  \item \solution{Mögliche Kategorien: Identitätsverwaltung, Datenhaltung, Ereignisverwaltung, Lastenverteilung, Monitoring.}
  \item \solution{Identitätsverwaltung für Login und Rechte. Datenhaltung für Produkte, Bestellungen und Kundendaten. Ereignisverwaltung für Meldungen wie \glqq Bestellung erstellt\grqq{}. Lastenverteilung für viele gleichzeitige Anfragen. Monitoring für Logs, Fehlerraten und Reaktionszeiten.}
  \item \solution{Für die Benachrichtigung nach einer Bestellung ist besonders die Ereignisverwaltung passend.}
\end{enumerate}

\section*{8. Reflexion}
\begin{enumerate}
  \item \solution{Zu viele Verantwortungen in einer Komponente erhöhen Kopplung, Komplexität und Fehlerrisiko.}
  \item \solution{Weil zusätzliche Komponenten auch zusätzlichen Betriebs- und Abstimmungsaufwand erzeugen können.}
  \item \solution{Wichtiger ist, welche fachliche oder technische Verantwortung im konkreten System wirklich gebraucht wird.}
\end{enumerate}

\end{document}
